\documentclass[12pt, a4paper]{article}

\usepackage{amsmath, amssymb}
\usepackage{fontspec}
\setmainfont{FreeSerif}
\usepackage{geometry}
\geometry{margin=1in}
\usepackage[colorlinks=true, linkcolor=blue, urlcolor=blue, citecolor=blue]{hyperref}

\setcounter{secnumdepth}{0} 

\begin{document}

\tableofcontents
\vspace{2em}
\hrule
\vspace{2em}

\section{Հիշեցում. Լագրանժի թեորեմը մեկ փոփոխականի ֆունկցիայի համար}

Եթե $f(x)$ ֆունկցիան՝

\begin{itemize}
    \item անընդհատ է $[a, b]$ հատվածում,
    \item դիֆերենցելի է $(a, b)$ միջակայքում,
\end{itemize}

ապա գոյություն ունի առնվազն մի $c \in (a, b)$ կետ այնպիսին, որ.

$$f'(c) = \frac{f(b) - f(a)}{b - a}$$

կամ այլ կերպ ձևակերպված՝ $f(b) - f(a) = f'(c)(b - a)$:

\vspace{1em}
\hrule
\vspace{1em}

\section{Լագրանժի բանաձևը շատ փոփոխականների ֆունկցիաների համար}

\textbf{Թեորեմ:} Դիցուք $f: D \to \mathbb{R}$, որտեղ $D \subset \mathbb{R}^m$ բաց տիրույթ է: Դիցուք $x_0, x \in D$ և նրանց միացնող $[x_0, x]$ հատվածը ամբողջությամբ պատկանում է $D$-ին ($[x_0, x] \subset D$):

Եթե $f$ ֆունկցիան դիֆերենցելի է $[x_0, x]$ հատվածում, ապա գոյություն ունի $\theta \in (0, 1)$ այնպիսին, որ.

$$f(x) - f(x_0) = \sum_{i=1}^{m} \frac{\partial f}{\partial x^i} (x_0 + \theta(x - x_0)) (x^i - x_0^i)$$

Այսինքն՝ ֆունկցիայի լրիվ աճը հավասար է մասնակի ածանցյալների և փոփոխականների աճերի արտադրյալների գումարին՝ հաշվարկված միջանկյալ $x_0 + \theta(x - x_0)$ կետում:

\textbf{Ապացույց}
Ապացուցման համար խնդիրը բերում ենք մեկ փոփոխականի դեպքին:

\textbf{Քայլ 1: Օժանդակ ֆունկցիայի ներմուծում}
Դիտարկենք հետևյալ մեկ փոփոխականի $F(t)$ ֆունկցիան, որտեղ $t \in [0, 1]$.

$$F(t) = f(t x_0 + (1 - t) x)$$

*(Նշում: Oգտագործված է նաև $F(t) = f(x + t(x_0 - x))$ տարբերակը, ինչը էապես նույնն է և ներկայացնում է հատվածի պարամետրական տեսքը:)*

\textbf{Քայլ 2: Ածանցում ըստ բարդ ֆունկցիայի կանոնի}
Քանի որ $f$-ը դիֆերենցելի է, իսկ հատվածի կոորդինատները գծային (հետևաբար դիֆերենցելի) ֆունկցիաներ են $t$-ից, ապա $F(t)$-ն նույնպես դիֆերենցելի է: Օգտվելով բարդ ֆունկցիայի ածանցման կանոնից (շղթայական կանոն).

$$\frac{dF}{dt} = \sum_{i=1}^{m} \frac{\partial f}{\partial x^i} (t x_0 + (1 - t) x) \cdot \frac{d}{dt} (t x_0^i + (1 - t) x^i)$$

Քանի որ $\frac{d}{dt} (t x_0^i + (1 - t) x^i) = x_0^i - x^i$, ստանում ենք.

$$\frac{dF}{dt} = \sum_{i=1}^{m} \frac{\partial f}{\partial x^i} (\dots) (x_0^i - x^i)$$

\textbf{Քայլ 3: Լագրանժի թեորեմի կիրառում $F(t)$-ի համար}
Կիրառելով մեկ փոփոխականի Լագրանժի թեորեմը $[0, 1]$ հատվածում $F(t)$ ֆունկցիայի համար՝ գոյություն ունի $\theta \in (0, 1)$, այնպես որ.

$$F(1) - F(0) = F'(\theta) (1 - 0) = F'(\theta)$$

Հաշվի առնելով, որ.

\begin{itemize}
    \item $F(1) = f(1 \cdot x_0 + 0 \cdot x) = f(x_0)$
    \item $F(0) = f(0 \cdot x_0 + 1 \cdot x) = f(x)$
\end{itemize}

Ստանում ենք.

$$f(x_0) - f(x) = \sum_{i=1}^{m} \frac{\partial f}{\partial x^i} (\theta x_0 + (1 - \theta) x) (x_0^i - x^i)$$

Բազմապատկելով $-1$-ով՝ կստանանք թեորեմի ձևակերպման մեջ նշված վերջնական տեսքը: \textbf{\#}

\vspace{1em}
\hrule
\vspace{1em}

\subsection{Պնդում}

Եթե $f: D \to \mathbb{R}$ ֆունկցիան դիֆերենցելի է կապակցված $D \subset \mathbb{R}^m$ բաց տիրույթում և նրա բոլոր մասնակի ածանցյալները զրո են տիրույթի բոլոր կետերում՝

$$\frac{\partial f}{\partial x^i} = 0, \quad \text{բոլոր } i=1, \dots, m \text{ համար,}$$

ապա $f(x) = C$ (հաստատուն) $D$ տիրույթում։

\textbf{Ապացույց (ուղղագիծ հատվածի համար)}
\begin{enumerate}
    \item \textbf{Կետերի ընտրություն:} Վերցնենք տիրույթի կամայական երկու $x_0$ և $x$ կետեր այնպես, որ նրանց միացնող հատվածը $[x_0, x]$ ամբողջությամբ ընկած լինի $D$ տիրույթի մեջ։
    \item \textbf{Լագրանժի բանաձևի կիրառում:} Ըստ Լագրանժի վերջավոր աճերի բանաձևի՝ գոյություն ունի միջանկյալ $\xi = x_0 + \theta(x - x_0)$ կետ ($\theta \in (0, 1)$), որի համար ճիշտ է հետևյալ հավասարությունը.
\end{enumerate}

$$f(x) - f(x_0) = \frac{\partial f}{\partial x^1}(\xi)(x^1 - x_0^1) + \frac{\partial f}{\partial x^2}(\xi)(x^2 - x_0^2) + \dots + \frac{\partial f}{\partial x^m}(\xi)(x^m - x_0^m)$$

\begin{enumerate}
    \setcounter{enumi}{2}
    \item \textbf{Ածանցյալների զրոյական արժեքը:} Քանի որ ըստ պայմանի բոլոր մասնակի ածանցյալները զրո են տիրույթի \textbf{ցանկացած} կետում, ապա դրանք զրո են նաև միջանկյալ $\xi$ կետում.
\end{enumerate}

$$\frac{\partial f}{\partial x^i}(\xi) = 0, \quad \text{բոլոր } i \text{-երի համար}$$

\begin{enumerate}
    \setcounter{enumi}{3}
    \item Տեղադրելով զրոները բանաձևի մեջ՝ ստանում ենք.
\end{enumerate}

$$f(x) - f(x_0) = 0 \cdot (x^1 - x_0^1) + 0 + \dots + 0 = 0$$

Հետևաբար՝ \textbf{$f(x) = f(x_0)$}։

\vspace{1em}
\hrule
\vspace{1em}

\subsubsection{Ընդհանուր դեպքը (կապակցված տիրույթ)}

Եթե $x$ և $x_0$ կետերը հեռու են և նրանց միացնող ուղիղ հատվածը դուրս է գալիս տիրույթի սահմաններից, մենք օգտվում ենք տիրույթի \textbf{կապակցվածության} հատկությունից։ Մենք կարող ենք $x_0$ և $x$ կետերը միացնել բեկյալով, որի յուրաքանչյուր հատված ընկած է $D$-ի մեջ։

Կիրառելով վերը նշված տրամաբանությունը բեկյալի յուրաքանչյուր հատվածի համար՝ կստանանք, որ ֆունկցիայի արժեքը չի փոխվում մի հատվածի սկզբից մինչև վերջ։ Արդյունքում՝ $f(x) = f(x_0)$ ամբողջ տիրույթում։

Մաթեմատիկական անալիզում «բաց կապակցված» բազմությունը (հաճախ անվանում են պարզապես \textbf{տիրույթ}) այն միջավայրն է, որտեղ մենք սովորաբար աշխատում ենք ֆունկցիաների հետ, որպեսզի խուսափենք «թռիչքներից» կամ մեկուսացված կետերից։

\textbf{Ի՞նչ է նշանակում «Բաց»}

Բազմությունը կոչվում է բաց, եթե նրա յուրաքանչյուր կետ \textbf{ներքին կետ} է։
Այսինքն՝ եթե վերցնես այդ բազմության ցանկացած $x$ կետ, դու միշտ կարող ես գտնել մի փոքր շրջակայք (գունդ), որն ամբողջությամբ տեղավորվում է այդ բազմության մեջ։

\begin{itemize}
    \item \textbf{Պարզ ասած.} Բաց բազմությունը չի պարունակում իր եզրագիծը (սահմանը)։
    \item \textbf{Օրինակ.} Մեկ չափողականության մեջ $(a, b)$ միջակայքը բաց է, իսկ $[a, b]$ հատվածը՝ փակ (քանի որ $a$ և $b$ ծայրակետերից դեպի դուրս քայլ անելիս դուրս ես գալիս բազմությունից)։
\end{itemize}

\vspace{1em}
\hrule
\vspace{1em}

\textbf{Ի՞նչ է նշանակում «Կապակցված»}

Տիրույթների դեպքում (բաց բազմությունների համար) կապակցվածությունը նշանակում է \textbf{գծային կապակցվածություն}։
Բազմությունը կոչվում է կապակցված, եթե նրա ցանկացած երկու կետ կարելի է միացնել այդ բազմությանը պատկանող \textbf{անընդհատ գծով} (կամ բեկյալով)։

\begin{itemize}
    \item \textbf{Պարզ ասած.} Բազմությունը «մեկ կտորից» է։ Այն բաղկացած չէ իրարից անջատ կղզյակներից։
    \item \textbf{Օրինակ.} Երկու առանձին շրջաններից բաղկացած բազմությունը կապակցված չէ, որովհետև մի շրջանից մյուսը գնալու համար պետք է «թռչել» դատարկ տարածության վրայով։
\end{itemize}

\textbf{Ինչո՞ւ է սա կարևոր նախորդ պնդման համար}

Երբ մենք ասում ենք, որ «եթե ածանցյալները զրո են, ապա ֆունկցիան հաստատուն է», \textbf{կապակցվածությունը վճռորոշ է}։

Եթե բազմությունը կապակցված չլիներ (օրինակ՝ բաղկացած լիներ երկու իրարից հեռու շրջաններից), ապա ֆունկցիան կարող էր մի շրջանի վրա լինել $f(x)=5$, իսկ մյուսի վրա՝ $f(x)=10$։ Երկու դեպքում էլ ածանցյալները կլինեին զրո, բայց ֆունկցիան ամբողջ բազմության վրա հաստատուն չէր լինի (քանի որ տարբեր կտորների վրա տարբեր արժեքներ ունի)։

\vspace{1em}
\hrule
\vspace{1em}

Հիմնվելով քո տրամադրած լուսանկարների վրա՝ ահա դասախոսության ամբողջական վերականգնումը, որն ընդգրկում է \textbf{ածանցյալն ըստ ուղղության} և \textbf{գրադիենտի} թեմաները:

\vspace{1em}
\hrule
\vspace{1em}

\section{Ածանցյալն ըստ ուղղության (Directional Derivative)}

\textbf{Սահմանում:} Դիցուք $f: D \to \mathbb{R}$, որտեղ $D \subset \mathbb{R}^m$ բաց տիրույթ է: $V$-ն միավոր վեկտոր է $\mathbb{R}^m$-ում ($|V|=1$):
Եթե գոյություն ունի հետևյալ վերջավոր սահմանը.

$$\lim_{\Delta x \to 0} \frac{f(x_0 + \Delta x \cdot V) - f(x_0)}{\Delta x}$$

ապա այն կոչվում է $f$ ֆունկցիայի \textbf{ածանցյալ $V$ վեկտորի ուղղությամբ} $x_0$ կետում և նշանակվում է $\frac{\partial f}{\partial V}(x_0)$:

$$\frac{\partial f}{\partial V}(x_0) = \lim_{\Delta x \to 0} \frac{f(x_0 + \Delta x \cdot V) - f(x_0)}{\Delta x} \quad (1) \text{}$$

\subsection{Մասնավոր դեպքեր (R\textasciicircum 2-ում):}

\begin{itemize}
    \item \textbf{$OX$ առանցքի ուղղությամբ:} Եթե $V = (1, 0)$, ապա սահմանը վերածվում է $x$ ըստ մասնակի ածանցյալի:
\end{itemize}

$$\frac{\partial f}{\partial (1,0)}(x_0, y_0) = \lim_{\Delta x \to 0} \frac{f(x_0 + \Delta x, y_0) - f(x_0, y_0)}{\Delta x} = \frac{\partial f}{\partial x} \text{}$$

\begin{itemize}
    \item \textbf{$OY$ առանցքի ուղղությամբ:} Եթե $V = (0, 1)$, ստանում ենք մասնակի ածանցյալն ըստ $y$-ի:
\end{itemize}

$$\frac{\partial f}{\partial (0,1)}(x_0, y_0) = \lim_{\Delta y \to 0} \frac{f(x_0, y_0 + \Delta y) - f(x_0, y_0)}{\Delta y} = \frac{\partial f}{\partial y} \text{}$$

\vspace{1em}
\hrule
\vspace{1em}

\subsection{Թեորեմ. Մասնակի ածանցյալների կապը ուղղությամբ ածանցյալի հետ}

\textbf{Թեորեմ:} Եթե $f$ ֆունկցիան դիֆերենցելի է $x_0 \in D$ կետում, ապա նրա ածանցյալը ցանկացած $V = (V^1, \dots, V^m)$ ուղղությամբ հաշվվում է հետևյալ բանաձևով:

$$\frac{\partial f}{\partial V}(x_0) = \frac{\partial f}{\partial x^1}(x_0)V^1 + \dots + \frac{\partial f}{\partial x^m}(x_0)V^m \quad (2) \text{}$$

\subsubsection{Ապացույց:}

Դիտարկենք օժանդակ $F(t) = f(x_0 + tV)$ ֆունկցիան: Քանի որ $f$-ը դիֆերենցելի է, իսկ կոորդինատային ֆունկցիաները գծային են, ըստ բարդ ֆունկցիայի ածանցման կանոնի:

$$\frac{\partial F}{\partial t} = \frac{\partial f}{\partial x^1}(x_0 + tV) \cdot V^1 + \dots + \frac{\partial f}{\partial x^m}(x_0 + tV) \cdot V^m \text{}$$

Տեղադրելով $t=0$, ստանում ենք:

$$\frac{\partial F}{\partial t}(0) = \frac{\partial f}{\partial x^1}(x_0)V^1 + \dots + \frac{\partial f}{\partial x^m}(x_0)V^m \text{}$$

Մյուս կողմից, ըստ ածանցյալի սահմանման:

$$\frac{\partial F}{\partial t}(0) = \lim_{\Delta x \to 0} \frac{F(0 + \Delta x) - F(0)}{\Delta x} = \frac{\partial f}{\partial V}(x_0) \text{}$$

Հետևաբար $(1) = (2)$, ինչը և պահանջվում էր ապացուցել:

\vspace{1em}
\hrule
\vspace{1em}

\subsection{Գրադիենտ (Gradient)}

\textbf{Սահմանում:} $x_0 \in D$ կետում $f$-ի մասնակի ածանցյալներից կազմված վեկտորը կոչվում է \textbf{գրադիենտ}:

$$\nabla f(x_0) = \left( \frac{\partial f}{\partial x^1}(x_0), \dots, \frac{\partial f}{\partial x^m}(x_0) \right) \text{}$$

\textbf{Կապը սկալյար արտադրյալի հետ:} Ուղղությամբ ածանցյալը հավասար է գրադիենտի և $V$ վեկտորի էվկլիդյան սկալյար արտադրյալին:

$$\frac{\partial f}{\partial V}(x_0) = \langle \nabla f(x_0), V \rangle \text{}$$

\vspace{1em}
\hrule
\vspace{1em}

\subsection{Աճի ամենաարագ ուղղությունը}

$\mathbb{R}^3$-ի դեպքում, որտեղ $V = (V^1, V^2, V^3)$, սկալյար արտադրյալը կարելի է գրել որպես:

$$\frac{\partial f}{\partial V}(x_0) = \|\nabla f(x_0)\| \cdot \|V\| \cdot \cos(\nabla f, V) \text{}$$

Քանի որ $\|V\| = 1$, ապա արժեքը կախված է միայն անկյունից:

\textbf{Առավելագույն աճ:} Քանի որ $\cos(\theta)$-ի առավելագույն արժեքը 1 է, սկալյար արտադրյալը (այսինքն՝ ուղղությամբ ածանցյալը) կլինի ամենամեծը, երբ $\theta = 0$: Իսկ $0$ աստիճանի անկյուն նշանակում է, որ $V$ վեկտորը և $\nabla f$ գրադիենտը նույն ուղղությամբ են:

\textbf{Զրոյական աճ:} Եթե վեկտորները ուղղահայաց են ($\theta = 90^\circ$), ապա $\cos(90^\circ) = 0$, և սկալյար արտադրյալը դառնում է $0$: Սա նշանակում է, որ այդ ուղղությամբ ֆունկցիան չի փոխվում (շարժվում ենք մակարդակի գծի երկայնքով):

\textbf{Նվազում:} Եթե վեկտորները հակառակ ուղղությամբ են ($\theta = 180^\circ$), ապա $\cos(180^\circ) = -1$, և մենք ունենում ենք ամենաարագ նվազումը:

\end{document}